\documentclass[letterpaper,11pt]{article}

% Shared LaTeX preamble for CV and Cover Letter

\usepackage{iftex}

\ifPDFTeX
  \usepackage[T2A]{fontenc}
  \usepackage[utf8]{inputenc}
\else
  \usepackage{fontspec}
\fi

\usepackage[english, russian]{babel}

\usepackage{latexsym}
\usepackage[empty]{fullpage}
\usepackage{titlesec}
\usepackage{marvosym}
\usepackage[usenames,dvipsnames]{color}
\usepackage{verbatim}
\usepackage{enumitem}
\usepackage[hidelinks]{hyperref}
\usepackage{fancyhdr}
\usepackage{tabularx}
\usepackage{graphicx}

\input{glyphtounicode}
\pdfgentounicode=1

\pagestyle{fancy}
\fancyhf{} % clear all header and footer fields
\fancyfoot{}
\renewcommand{\headrulewidth}{0pt}
\renewcommand{\footrulewidth}{0pt}

\urlstyle{same}
\raggedbottom
\raggedright

\newcommand{\plusplus}{\texttt{++}}
\newcommand{\cpp}{C\plusplus}

% Personal profile and contact information (Single Source of Truth)

\newcommand{\fullName}{Горобченко Александр}
\newcommand{\phoneFormatted}{+7 (995) 375-50-07}
\newcommand{\phoneUrl}{+79953755007}
\newcommand{\emailAddr}{alexo375@yandex.ru}
\newcommand{\telegramUser}{alexo375}
\newcommand{\githubUser}{alexo375}
\newcommand{\githubUrl}{https://github.com/\githubUser}
\newcommand{\birthDate}{11.09.2004}
\newcommand{\locationInfo}{Рязань / Готов к релокации в Москву}
\newcommand{\photoFile}{me.png}

% Compatibility aliases
\providecommand{\username}{\githubUser}
\providecommand{\github}{\githubUrl}

% Reusable Header Macro
\newcommand{\renderheader}[1]{%
  \begin{center}
    \begin{minipage}[b]{0.80\textwidth}
        \textbf{\Huge \scshape \fullName} \\ \vspace{4pt}
        \large #1 \\ \vspace{4pt}
        \small \href{tel:\phoneUrl}{\phoneFormatted} $|$ 
        \href{mailto:\emailAddr}{\underline{\emailAddr}} $|$ 
        \href{https://t.me/\telegramUser}{\underline{t.me/\telegramUser}} $|$
        \href{\githubUrl}{\underline{github.com/\githubUser}} \\ \vspace{2pt}
        \small \birthDate \\ \vspace{2pt}
        \small \locationInfo
    \end{minipage}%
    \begin{minipage}[b]{0.18\textwidth}
        \raggedleft
        \includegraphics[width=4cm, keepaspectratio]{\photoFile} 
    \end{minipage}
  \end{center}
}


% Adjust margins
\addtolength{\oddsidemargin}{-0.5in}
\addtolength{\evensidemargin}{-0.5in}
\addtolength{\textwidth}{1in}
\addtolength{\topmargin}{-.5in}
\addtolength{\textheight}{1.0in}

% Настройки абзацев для сопроводительного письма (без красной строки, с отступом между абзацами)
\setlength{\parindent}{0pt}
\setlength{\parskip}{1em}

\begin{document}

%----------ЗАГОЛОВОК----------
\renderheader{Data Science \& Machine Learning}

\vspace{0.5cm}
\hrule
\vspace{0.5cm}

\textbf{Тема:} Отклик на позицию в области ML / DS в Ozon

Здравствуйте, команда Ozon!

Меня интересует развитие в командах машинного обучения и анализа данных Ozon. Задачи построения и внедрения ML-моделей на больших данных, их оптимизации и доведения до продакшена вызывают у меня сильный исследовательский и инженерный интерес.

О себе и релевантном опыте:
\begin{itemize}[leftmargin=0.2in, itemsep=2pt, topsep=0pt]
	\item Обучаюсь на 1 курсе магистратуры Центрального Университета по направлению <<Математика и компьютерные науки>>, ML-трек. Имею крепкую математическую и алгоритмическую базу и готов работать 40 часов в неделю.
	\item Имею практический опыт полного цикла работы с ML и Deep Learning: от сборки сырых данных и разметки датасетов, проектирования признаков до обучения нейросетевых архитектур и классических моделей, с упором на качество и оптимизацию решений.
\end{itemize}

Мне нравится открытый и инженерный подход Ozon Tech, и я хочу применить свои навыки в работе с большими данными и ML-моделями для развития продуктов и сервисов Ozon, чтобы ежедневно приносить пользу пользователям.

Буду рад познакомиться на интервью и обсудить возможности.

С уважением,\\
\fullName

\end{document}
